\documentclass[a4paper,11pt]{article}
\usepackage{exam-style-341}
\examinehead{341 农业综合知识三 · 数据库分章练习}{第02章\quad 关系模型与关系代数\quad 参考答案}
\begin{document}

\unitanswerscover{341 农业综合知识三 · 数据库分章练习}{第02章\quad 关系模型与关系代数 参考答案}{答案与解析按练习册题号顺序编排；判断题先给出判断，再给出改正后的正确表述；综合题给出关系代数表达式的逐步构造与属性集闭包运算的完整推导。}

\answerpart{一、单项选择题（10 题，每题 2 分，共 20 分）}

\q \textbf{B}。\begin{analysis}在关系模型中，二维表称为关系，表中的一行称为元组（记录），一列称为属性（字段），属性的取值范围称为域。\end{analysis}

\q \textbf{B}。\begin{analysis}候选码是能够唯一标识一个元组且不含多余属性的属性组，即极小的超码。超码可能含有多余属性；主码是从候选码中选定的一个；外码是参照另一关系主码的属性组。题干强调“不含多余属性”，故选候选码。\end{analysis}

\q \textbf{A}。\begin{analysis}包含在任何一个候选码中的属性称为主属性；不包含在任何一个候选码中的属性称为非主属性（非码属性）。\end{analysis}

\q \textbf{A}。\begin{analysis}实体完整性规则：若属性（组）$A$ 是基本关系 $R$ 的主属性，则 $A$ 不能取空值。参照完整性针对外码；用户定义完整性是用户针对具体应用定义的约束；域完整性指属性取值必须来自其值域，不是三类完整性中的标准名称。\end{analysis}

\q \textbf{A}。\begin{analysis}参照完整性规则：若属性（组）$F$ 是基本关系 $R$ 的外码，它与基本关系 $S$ 的主码 $K$ 相对应，则 $R$ 中每个元组在 $F$ 上的值必须为：或者取空值，或者等于 $S$ 中某个元组的主码值。B、C 都只取了一种情形，不全面；D 与规则相矛盾。\end{analysis}

\q \textbf{B}。\begin{analysis}$\sigma$ 是选择运算，从关系中选出满足给定条件的元组（行）；$\pi$ 是投影运算，从关系中选出若干属性列；$\bowtie$ 是连接运算；$\div$ 是除运算。\end{analysis}

\q \textbf{B}。\begin{analysis}$R-(R-S)$ 中，$R-S$ 表示属于 $R$ 而不属于 $S$ 的元组，再从 $R$ 中减去这部分，剩下的正是既属于 $R$ 又属于 $S$ 的元组，即 $R\cap S$。这是交运算的等价表达式。\end{analysis}

\q \textbf{B}。\begin{analysis}自然连接要求两个关系中进行比较的分量必须是同名属性组，并在结果中去掉重复的属性列。$R(A,B)$ 与 $S(B,C)$ 的公共属性为 $B$，结果关系模式为 $(A,B,C)$，共 3 个属性。\end{analysis}

\q \textbf{C}。\begin{analysis}关系代数的五种基本运算是并、差、笛卡尔积、选择和投影。交、连接、除等运算都可以用这五种基本运算表达，属于导出运算（导出运算），因此交不属于基本运算。\end{analysis}

\q \textbf{A}。\begin{analysis}关系的性质包括：每个属性都是不可再分的原子值；元组的次序可以任意交换；属性的次序可以任意交换；不允许出现两个完全相同的元组。B 违反“不允许重复元组”；C 违反属性的原子性；D 违反“属性次序无关”。故选 A。\end{analysis}

\answerpart{二、多项选择题（5 题，每题 4 分，共 20 分）}

\q \textbf{ABCD}。\begin{analysis}关系的性质包括：属性不可再分（原子性）；元组次序无关；属性次序无关；不允许出现完全相同的元组。四项均正确。\end{analysis}

\q \textbf{ABC}。\begin{analysis}关系数据库的三类完整性约束是实体完整性、参照完整性和用户定义完整性。D 数据独立性是数据库系统的特点，不属于完整性约束。\end{analysis}

\q \textbf{ABCD}。\begin{analysis}五种基本运算是并（$\cup$）、差（$-$）、笛卡尔积（$\times$）、选择（$\sigma$）和投影（$\pi$）。四个选项均为基本运算。\end{analysis}

\q \textbf{ABC}。\begin{analysis}外码是本关系中的一个属性组，它不是本关系的主码，但参照另一个关系的主码；外码的取值可以为空值（当外码不是主属性且应用语义允许时），也可以等于被参照关系中某个元组的主码值。D 错误：外码不是本关系的主码，否则就失去“参照”的意义。\end{analysis}

\q \textbf{ABC}。\begin{analysis}A、C 都是交运算的等价表达式，正确；B 并运算满足交换律，正确。D 错误：差运算不满足交换律，一般情况下 $R-S\ne S-R$。\end{analysis}

\answerpart{三、判断题（5 题，每题 2 分，共 10 分）}

\q ×。改正：只有满足关系基本性质的二维表才是关系，即每个属性不可再分、元组不重复、行序和列序无关等；结构不符合这些要求的二维表不是关系。\begin{analysis}关系模型对关系有明确限定：属性必须是原子值，不允许表中有表；不允许出现完全相同的元组；元组和属性的次序都无关。因此“任何二维表都是关系”的说法不成立。\end{analysis}

\q ×。改正：一个关系可以有多个候选码，设计时从中选定一个作为主码。\begin{analysis}候选码是极小超码，能唯一标识元组的属性组可能不止一个。例如学生关系中“学号”和“身份证号”都可以唯一标识学生，二者都是候选码；主码只是其中被选定的一个，一个关系只有一个主码。\end{analysis}

\q √。\begin{analysis}参照完整性规则允许外码取空值（表示“尚未确定”或“没有参照”），也允许外码等于被参照关系中某个元组的主码值。但若外码同时是本关系的主属性，则受实体完整性约束不能取空值。\end{analysis}

\q √。\begin{analysis}笛卡尔积 $R\times S$ 的每个元组由 $R$ 的一个元组与 $S$ 的一个元组拼接而成，因此元组数为 $|R|\times|S|=m\times n$，属性个数为 $R$ 与 $S$ 的属性个数之和。\end{analysis}

\q √。\begin{analysis}关系是元组的集合，集合中不允许重复元素；投影后若出现重复元组，必须去掉重复元组，这是投影运算与一般“选列”操作的重要区别。\end{analysis}

\answerpart{四、填空题（5 题，每题 2 分，共 10 分）}

\q 第一空：\textbf{选择}；第二空：\textbf{投影}。\begin{analysis}五种基本关系代数运算是并、差、笛卡尔积、选择和投影；交、连接、除等都可以由它们导出。\end{analysis}

\q 第一空：\textbf{选择}；第二空：\textbf{投影}；第三空：\textbf{连接}。\begin{analysis}$\sigma$ 为选择，$\pi$ 为投影，$\bowtie$ 为连接（最常用的是自然连接）。\end{analysis}

\q 第一空：\textbf{$ABCD$}；第二空：\textbf{$A$}。\begin{analysis}$A^{+}$：初始为 $\{A\}$；由 $A\to B$ 得 $\{A,B\}$；由 $B\to C$ 得 $\{A,B,C\}$；由 $C\to D$ 得 $\{A,B,C,D\}$，故 $A^{+}=ABCD$。$A$ 能推出全部属性，且 $A$ 是单属性、不存在更小的真子集，故 $A$ 是候选码（也是唯一的候选码）。\end{analysis}

\q 第一空：\textbf{$AB$}；第二空：\textbf{$A$、$B$}；第三空：\textbf{$C$}。\begin{analysis}$F=\{AB\to C\}$ 中只有 $AB$ 能推出 $C$，$A^{+}=A$、$B^{+}=B$ 都不能推出全部属性，故候选码为 $AB$。包含在候选码中的属性 $A$、$B$ 为主属性，$C$ 为非主属性。\end{analysis}

\q 第一空：\textbf{元组}；第二空：\textbf{属性}；第三空：\textbf{域}。\begin{analysis}关系模型用二维表表示关系：一行是一个元组，一列是一个属性，属性的取值范围称为域（值域）。\end{analysis}

\answerpart{五、简答题（2 题，每题 8 分，共 16 分）}

\q 参考答案要点：

\textbf{（1）超码（2 分）}：能够唯一标识关系中一个元组的属性组，可能含有多余属性。例如在学生关系中，（学号，姓名）能唯一标识学生，是超码，但“姓名”是多余的。

\textbf{（2）候选码（2 分）}：能够唯一标识一个元组且不含多余属性的属性组，即极小的超码。例如“学号”和“身份证号”都能唯一标识学生且本身不能再约简，二者都是候选码。一个关系可以有多个候选码。

\textbf{（3）主码（2 分）}：从多个候选码中选定一个作为主码，用于唯一标识元组；一个关系只有一个主码。例如选定“学号”为主码。

\textbf{（4）外码（2 分）}：本关系中的一个属性组，它不是本关系的主码，但参照另一个关系的主码。例如选课关系 SC 中的“学号”参照学生关系 Student 的主码“学号”，SC 的“学号”就是外码。

\textbf{（5）联系}：候选码是极小超码，主码是选定的候选码，三者都用于唯一标识元组；外码用于表示关系之间的联系，其取值必须符合参照完整性。超码 $\supseteq$ 候选码 $\supseteq$ 主码，外码与它们属于不同关系的概念。
\begin{analysis}评分要点：四个概念的定义各 1.5 分，共 6 分；能说明四者的联系与区别 2 分。常见错误：把“超码”与“候选码”混为一谈，或认为主码可以有多个、候选码只能有一个。\end{analysis}

\q 参考答案要点：

\textbf{（1）实体完整性（3 分）}
规则内容：若属性（组）$A$ 是基本关系 $R$ 的主属性，则 $A$ 不能取空值。
说明：主码唯一标识元组，若主属性取空值，就无法区分元组。违反时一般拒绝插入或修改。

\textbf{（2）参照完整性（3 分）}
规则内容：若属性（组）$F$ 是基本关系 $R$ 的外码，它与基本关系 $S$ 的主码 $K$ 相对应，则 $R$ 中每个元组在 $F$ 上的值必须为：或者取空值，或者等于 $S$ 中某个元组的主码值。
说明：外码取空值表示“尚未确定”；取非空值时必须能在被参照关系中找到对应元组。违反时一般拒绝执行，也可按应用语义采用级联删除或置空。

\textbf{（3）用户定义完整性（2 分）}
规则内容：用户针对具体应用环境定义的约束条件，如“成绩在 0 到 100 之间”“性别只能取男或女”。
说明：这类约束由 DBMS 提供机制、由用户定义，违反时拒绝执行。
\begin{analysis}评分要点：三类完整性的名称（各 0.5 分，共 1.5 分）、规则内容表述（各 1.5 分左右）、违反时的处理（各 1 分）。常见错误：把实体完整性说成“主码唯一”，或漏掉参照完整性中“或者取空值”的情形。\end{analysis}

\answerpart{六、综合题（2 题，每题 12 分，共 24 分）}

\q 参考答案：

\textbf{（1）查询计算机系年龄小于 20 岁的男学生姓名（3 分）}
\[
\pi_{\text{Sname}}\left(\sigma_{\text{Sdept}='CS'\ \wedge\ \text{Ssex}='男'\ \wedge\ \text{Sage}<20}(\text{Student})\right)
\]
先对 Student 做选择，同时用三个条件筛选系别、性别和年龄，再对结果投影出 Sname。

\textbf{（2）查询选修了课程号 'C01' 的课程的学生姓名（3 分）}
\[
\pi_{\text{Sname}}\left(\text{Student}\bowtie\sigma_{\text{Cno}='C01'}(\text{SC})\right)
\]
先在 SC 中选出课程号为 C01 的元组，再与学生关系做自然连接（公共属性 Sno），最后投影出姓名。也可写成
\[
\pi_{\text{Sname}}\left(\sigma_{\text{Cno}='C01'}(\text{Student}\bowtie\text{SC})\right),
\]
两种写法等价。

\textbf{（3）查询没有选修任何课程的学生的学号（3 分）}
\[
\pi_{\text{Sno}}(\text{Student})-\pi_{\text{Sno}}(\text{SC})
\]
先分别求出全部学生的学号和选过课的学生的学号，再用差运算求出“在全体学生中而不在选课学生中”的学号。

\textbf{（4）查询选修了全部课程的学生的学号（3 分）}
\[
\pi_{\text{Sno,Cno}}(\text{SC})\div\pi_{\text{Cno}}(\text{Course})
\]
除运算的含义是：设被除关系 $R$ 的属性集为 $X\cup Y$，除关系 $S$ 的属性集为 $Y$，则 $R\div S$ 的结果属性集为 $X$，其中的元组 $x$ 满足“对 $S$ 中的每一个元组 $y$，$(x,y)$ 都出现在 $R$ 中”。本题中 $R=\pi_{\text{Sno,Cno}}(\text{SC})$，$S=\pi_{\text{Cno}}(\text{Course})$，结果为学号集合，表示该学生选修了 Course 中出现的每一门课程。
\begin{analysis}评分要点：每题 3 分，其中关系代数的运算选择正确 2 分、表达式书写完整（选择条件、连接属性、投影属性齐全）1 分。易错点：（1）漏写连接条件或把选择条件写错；（2）用笛卡尔积而不加连接条件；（3）第（3）题误用“选课表中没有出现的学号”而漏掉“全体学生”这一来源；（4）第（4）题不会用除运算，或把被除关系误写成整个 SC（含 Grade，导致除运算属性不匹配）。\end{analysis}

\q 参考答案：

\textbf{（1）单属性闭包（3 分）}
\begin{itemize}[leftmargin=1.6em, itemsep=1pt]
  \item $A^{+}$：初始 $\{A\}$；由 $A\to B$ 得 $\{A,B\}$；$B\to A$ 未产生新属性，故 $A^{+}=\{A,B\}$。
  \item $B^{+}$：初始 $\{B\}$；由 $B\to A$ 得 $\{A,B\}$；$A\to B$ 未产生新属性，故 $B^{+}=\{A,B\}$。
  \item $C^{+}$：初始 $\{C\}$；$F$ 中左部含 $C$ 的只有 $AC\to D$，而 $A$ 不在集合中，不能使用，故 $C^{+}=\{C\}$。
\end{itemize}

\textbf{（2）属性集闭包与全部候选码（4 分）}
\begin{itemize}[leftmargin=1.6em, itemsep=1pt]
  \item $AC^{+}$：初始 $\{A,C\}$；由 $A\to B$ 得 $\{A,B,C\}$；由 $AC\to D$ 得 $\{A,B,C,D\}$，故 $AC^{+}=ABCD$。
  \item $BC^{+}$：初始 $\{B,C\}$；由 $B\to A$ 得 $\{A,B,C\}$；由 $AC\to D$ 得 $\{A,B,C,D\}$，故 $BC^{+}=ABCD$。
  \item 由于 $AC^{+}$ 和 $BC^{+}$ 都包含全部属性，且 $A$、$B$、$C$ 单独都不能推出全部属性，故 $AC$ 与 $BC$ 都是候选码。
  \item 进一步说明没有其他候选码：属性 $C$ 不会出现在任何函数依赖的右部，即 $C$ 不能由其他属性推出，所以每个候选码都必须包含 $C$；候选码中还必须含有能推出 $D$ 的 $A$ 或 $B$。含 $C$ 且含 $A$ 或 $B$ 的极小组合只有 $AC$ 和 $BC$，故全部候选码为 $AC$、$BC$。
\end{itemize}

\textbf{（3）主属性与非主属性（2 分）}
主属性：包含在候选码 $AC$ 或 $BC$ 中的属性，即 $A$、$B$、$C$；
非主属性：不包含在任何候选码中的属性，即 $D$。

\textbf{（4）候选码的极小性说明（3 分）}
$A$ 单独不能成为候选码：$A^{+}=AB\ne ABCD$，缺少 $C$ 和 $D$；$B$ 同理；$C$ 单独更不能成为候选码（$C^{+}=C$）。
对 $AC$：去掉 $A$ 后 $C^{+}=C$，不含全部属性；去掉 $C$ 后 $A^{+}=AB$，也不含全部属性。因此 $AC$ 的任何真子集都不能唯一标识元组，$AC$ 不含多余属性，且 $AC^{+}=ABCD$，故 $AC$ 是候选码。对 $BC$ 同理：去掉 $B$ 得 $C^{+}=C$，去掉 $C$ 得 $B^{+}=AB$，均不能推出全部属性，故 $BC$ 也是候选码。
\begin{analysis}评分要点：闭包运算每一步的推导过程（3 分）；$AC^{+}$、$BC^{+}$ 及候选码结论（4 分）；主属性、非主属性判断（2 分）；极小性论证（3 分）。常见错误：（1）求闭包时只写结果不写推导过程；（2）漏掉 $B\to A$ 导致 $BC^{+}$ 求错；（3）认为 $A$、$B$ 也是候选码（$A^{+}=AB$ 不含 $C$、$D$）；（4）把 $A$、$B$ 以外的属性都当成非主属性而漏掉 $C$。\end{analysis}

\end{document}
